% ---
% RESUMOS
% ---

% resumo em português
\setlength{\absparsep}{18pt} % ajusta o espaçamento dos parágrafos do resumo
\begin{resumo}
Este trabalho apresenta o desenvolvimento de uma Fechadura Biométrica Inteligente implementada na FPGA DE2-115 (Cyclone IV EP4CE115F29C7) como projeto integrador da disciplina ENGG52 da UFBA. O sistema realiza reconhecimento facial em tempo real através de uma Rede Neural Convolucional embarcada, classificando indivíduos entre 19 classes — 18 membros autorizados e uma classe nativa de rejeição — sem necessidade de limiar de confiança externo.
A solução cobre o pipeline completo de ponta a ponta: uma câmera conectada ao PC captura o rosto do usuário, que é detectado pelo algoritmo Haar Cascade, normalizado com CLAHE, redimensionado para 32×32 pixels e transmitido via UART à placa. Na FPGA, uma Máquina de Estados Finitos orquestra o pipeline de inferência composto pelos módulos de convolução 3×3 com ReLU (4 filtros), \textit{Max Pooling} 2×2, \textit{Flatten} e camada Densa 900×19, cujos pesos foram treinados em Python com Keras, quantizados para o formato de ponto fixo Q1.7 e armazenados em blocos de memória M9K. O resultado é exibido em um monitor VGA com o nome do indivíduo reconhecido e acionamento de LEDs indicando liberação ou negação de acesso.
As validações quantitativas realizadas ao longo do projeto confirmam a fidelidade da implementação em Verilog. O sistema integra os Artefatos 1 a 4 do cronograma proposto, abrangendo aquisição e treinamento off-line, arquitetura RTL, \textit{pipeline} de visão e comunicação, e integração final.


 \textbf{Palavras-chave}: FPGA. reconhecimento facial. CNN embarcada. quantização de ponto fixo. verilog. DE2-115.
\end{resumo}
